\documentclass[11pt,a4paper]{article}
\usepackage[utf8]{inputenc}
\usepackage[T1]{fontenc}
\usepackage[french]{babel}
\usepackage{geometry}
\geometry{margin=2.5cm}
\usepackage{graphicx}
\usepackage{xcolor}
\usepackage{listings}
\usepackage{hyperref}
\usepackage{fancyhdr}
\usepackage{enumitem}
\usepackage{tcolorbox}
\usepackage{tikz}
\usepackage{amssymb}


\usepackage{fontawesome5}

% Configuration des listes
\setlist[itemize]{leftmargin=*, itemsep=3pt, topsep=3pt}
\setlist[enumerate]{leftmargin=*, itemsep=3pt, topsep=3pt}

% Configuration du code
\lstset{
    basicstyle=\ttfamily\small,
    breaklines=true,
    frame=single,
    backgroundcolor=\color{gray!10},
    keywordstyle=\color{blue},
    commentstyle=\color{green!60!black},
    stringstyle=\color{orange},
    tabsize=4,
    showstringspaces=false
}

% En-tête et pied de page
\pagestyle{fancy}
\fancyhf{}
\fancyhead[L]{\small Guide de Travail - Chatbot SUP'PTIC}
\fancyhead[R]{\small Workflow Git}
\fancyfoot[C]{\thepage}

% Environnements personnalisés
\newtcolorbox{important}[1][]{
  colback=red!5!white,
  colframe=red!75!black,
  fonttitle=\bfseries,
  title=#1
}

\newtcolorbox{info}[1][]{
  colback=blue!5!white,
  colframe=blue!75!black,
  fonttitle=\bfseries,
  title=#1
}

\newtcolorbox{astuce}[1][]{
  colback=green!5!white,
  colframe=green!75!black,
  fonttitle=\bfseries,
  title=#1
}

\newtcolorbox{exemple}[1][]{
  colback=orange!5!white,
  colframe=orange!75!black,
  fonttitle=\bfseries,
  title=#1
}

\begin{document}

% Page de titre
\begin{titlepage}
    \centering
    \vspace*{2cm}
    {\Huge \textbf{Guide de Travail}} \\[0.5cm]
    {\LARGE \textbf{Workflow Git \& Organisation}}
    
    \vspace{1cm}
    
    {\Large Projet Chatbot SUP'PTIC}
    
    \vspace{1.5cm}
    
    \begin{minipage}{0.8\textwidth}
        \centering
        \small Comment travailler efficacement en équipe \\
        avec Git, les branches et les Pull Requests
    \end{minipage}
    
    \vfill
    
    \begin{center}
        \textbf{Club Informatique SUP'PTIC} \\
        Février 2026
    \end{center}
    
\end{titlepage}

\tableofcontents
\clearpage

% Section 1 : Organisation du Dépôt
\section{Organisation du Dépôt Git}
\label{sec:organisation}

\subsection{Structure des Branches}

Notre dépôt Git est organisé autour de \textbf{deux branches principales protégées} :

\begin{info}[Branches Principales]
\begin{itemize}
    \item \textbf{main} : Version stable et finale du projet
    \begin{itemize}
        \item Code testé et validé
        \item Prêt pour démonstration/déploiement
        \item \faIcon{lock} \textcolor{red}{PROTÉGÉE - Accès restreint}
    \end{itemize}
    
    \item \textbf{dev} : Branche de développement principal
    \begin{itemize}
        \item Intégration de toutes les features
        \item Tests et validation avant merge vers main
        \item \faIcon{lock} \textcolor{red}{PROTÉGÉE - Accès restreint}
    \end{itemize}
\end{itemize}
\end{info}

\subsection{Branches de Travail par Équipe}

Chaque équipe dispose de sa propre branche pour travailler :

\begin{itemize}
    \item \textbf{backend} : Équipe Backend (API Django, TF-IDF)
    \item \textbf{frontend} : Équipe Frontend (HTML/CSS/JS)
    \item \textbf{database} : Équipe Base de Données (Models, migrations)
    \item \textbf{data} : Équipe Structuration des Données (CSV, génération)
\end{itemize}

\begin{important}[Règle Fondamentale]
\textbf{Vous ne travaillez JAMAIS directement sur main ou dev.}

Vous créez et travaillez sur votre branche d'équipe, puis vous proposez vos changements via une Pull Request.
\end{important}

\clearpage

% Section 2 : Règles du Dépôt
\section{Règles du Dépôt }
\label{sec:regles}

\subsection{Règle \#1 : Ne touchez pas à main et dev directement}

\begin{itemize}
    \item[\faIcon{times}] Ces deux branches sont \textbf{protégées}
    \item[\faIcon{times}] Vous ne pouvez pas les supprimer
    \item[\faIcon{times}] Vous ne pouvez pas y envoyer du code directement (push)
\end{itemize}

\vspace{0.5cm}

\subsection{Règle \#2 : Travaillez sur vos propres branches}

\begin{itemize}
    \item[\faIcon{check}] Exemple : \texttt{backend}, \texttt{frontend}, \texttt{database}, \texttt{data}
    \item[\faIcon{check}] Faites vos modifications et commits là-bas
    \item[\faIcon{check}] Testez votre code avant de proposer
\end{itemize}

\vspace{0.5cm}

\subsection{Règle \#3 : Créez une Pull Request (PR) quand vous êtes prêts}

\begin{itemize}
    \item[\faIcon{code-branch}] Une PR = "Voilà mon travail, est-ce qu'on peut l'ajouter dans dev ?"
    \item[\faIcon{user-check}] Les chefs de projet vérifient et valident
    \item[\faIcon{check-circle}] Après validation, votre code est fusionné dans dev
\end{itemize}

\vspace{0.5cm}

\clearpage

% Section 3 : Workflow Résumé
\section{Workflow Résumé pour l'Équipe}
\label{sec:workflow}

\begin{enumerate}
    \item \textbf{Vous codez} sur vos branches personnelles (ex: \texttt{backend})
    
    \item \textbf{Vous proposez} vos changements via une Pull Request
    
    \item \textbf{Les chefs de projet} valident et fusionnent dans \texttt{dev}
    
    \item \textbf{Quand tout est stable}, \texttt{dev} est fusionné dans \texttt{main}
\end{enumerate}

\vspace{1cm}

\begin{figure}[h]
\centering
\begin{tikzpicture}[node distance=4.5cm, auto]
    % Styles
    \tikzstyle{branch} = [rectangle, rounded corners, minimum width=3cm, minimum height=1cm, text centered, draw=black, fill=blue!20]
    \tikzstyle{protected} = [rectangle, rounded corners, minimum width=3cm, minimum height=1cm, text centered, draw=black, fill=red!20]
    \tikzstyle{arrow} = [thick,->,>=stealth]
    
    % Nodes
    \node (team) [branch] {Branches équipes\newline(backend, frontend...)};
\node (dev) [protected, right of=team, xshift=2cm] {dev\newline(protégée)};
\node (main) [protected, right of=dev, xshift=2cm] {main\newline(protégée)};

    
    % Arrows
    \draw [arrow] (team) -- node[above] {Pull Request} (dev);
    \draw [arrow] (dev) -- node[above] {Validation} (main);
\end{tikzpicture}
\caption{Flux de travail Git simplifié}
\end{figure}

\clearpage

% Section 4 : Démarrage - Cloner le Dépôt
\section{Démarrer : Cloner le Dépôt}
\label{sec:demarrage}

\subsection{Étape 1 : Cloner le Dépôt sur Votre Machine}

\begin{lstlisting}[language=bash, caption=Dans le terminal VS Code ou PowerShell]
# Cloner le depot
git clone https://github.com/nkoumougrinnel/ChatBot.git

# Entrer dans le dossier
cd ChatBot
\end{lstlisting}

\subsection{Étape 2 : Vérifier les Branches Disponibles}

\begin{lstlisting}[language=bash]
# Lister toutes les branches (locales et distantes)
git branch -a
\end{lstlisting}

\textbf{Résultat attendu :}
\begin{lstlisting}
* main
  remotes/origin/backend
  remotes/origin/frontend
  remotes/origin/database
  remotes/origin/data
  remotes/origin/dev
  remotes/origin/main
\end{lstlisting}

\begin{astuce}[Comprendre les branches]
\begin{itemize}
    \item Les branches avec \texttt{remotes/origin/} sont sur GitHub (serveur distant)
    \item L'étoile \texttt{*} indique la branche actuelle
\end{itemize}
\end{astuce}

\clearpage

% Section 5 : Récupérer Votre Branche de Travail
\section{Récupérer Votre Branche de Travail}
\label{sec:recuperer}

\subsection{Cas 1 : Vous êtes dans l'Équipe Frontend}

\begin{lstlisting}[language=bash]
# Se positionner sur la branche frontend
git checkout frontend
\end{lstlisting}

\textbf{Que se passe-t-il ?}
\begin{itemize}
    \item[\faIcon{check}] Git crée automatiquement la branche \texttt{frontend} en local
    \item[\faIcon{check}] Git la synchronise avec \texttt{origin/frontend} (GitHub)
    \item[\faIcon{check}] Vous êtes maintenant sur la branche \texttt{frontend}
\end{itemize}

\subsection{Cas 2 : Vous êtes dans l'Équipe Backend}

\begin{lstlisting}[language=bash]
# Se positionner sur la branche backend
git checkout backend
\end{lstlisting}

\subsection{Mettre à Jour Votre Branche}

\textbf{Important :} Avant de commencer à coder, récupérez les dernières modifications :

\begin{lstlisting}[language=bash]
# Recuperer les derniers changements
git pull origin frontend
\end{lstlisting}

\textit{(Remplacez \texttt{frontend} par le nom de votre branche)}

\subsection{Vérifier Que Vous Êtes sur la Bonne Branche}

\begin{lstlisting}[language=bash]
# Verifier la branche actuelle
git branch
\end{lstlisting}

\textbf{Résultat attendu :}
\begin{lstlisting}
  backend
  database
* frontend    <- L'etoile indique votre branche actuelle
  main
\end{lstlisting}

\clearpage

% Section 6 : Travailler sur Votre Branche
\section{Travailler sur Votre Branche}
\label{sec:travailler}

\subsection{Organisation des Dossiers}

Quand vous ouvrez le projet dans VS Code, vous voyez tous les dossiers :

\begin{lstlisting}
ChatBot/
├── backend/        <- Backend et Database travaillent ICI
├── frontend/       <- Frontend travaille ICI
├── data/           <- Structuration et Database travaillent ICI
├── docs/
└── README.md
\end{lstlisting}

\begin{important}[Règle de Travail]
\textbf{Même si vous voyez tous les dossiers, vous devez travailler UNIQUEMENT dans le dossier de votre équipe.}

\vspace{0.3cm}
Exemples :
\begin{itemize}
    \item Équipe Frontend → Modifiez seulement les fichiers dans \texttt{frontend/}
    \item Équipe Backend → Modifiez seulement les fichiers dans \texttt{backend/}
    \item Équipe Database → Modifiez seulement les fichiers dans \texttt{backend/} (models, migrations)
    \item Équipe Data-struct → Modifiez seulement les fichiers dans \texttt{data/}
\end{itemize}
\end{important}

\subsection{Faire des Modifications}

\begin{enumerate}
    \item Ouvrez les fichiers de votre dossier
    \item Codez, testez, corrigez
    \item Sauvegardez vos fichiers (\texttt{Ctrl+S})
\end{enumerate}

\subsection{Enregistrer Vos Changements (Commit)}

\begin{lstlisting}[language=bash]
# Voir les fichiers modifies
git status

# Ajouter tous les fichiers modifies
git add .

# Creer un commit avec un message clair
git commit -m "Ajout interface chat + styles responsive"
\end{lstlisting}

\begin{astuce}[Messages de Commit Clairs]
\textbf{Mauvais exemples :}
\begin{itemize}
    \item[\faIcon{times}] "update"
    \item[\faIcon{times}] "fix"
    \item[\faIcon{times}] "changes"
\end{itemize}

\textbf{Bons exemples :}
\begin{itemize}
    \item[\faIcon{check}] "Ajout endpoint POST /api/chatbot/ask/"
    \item[\faIcon{check}] "Correction bug affichage feedback"
    \item[\faIcon{check}] "Import 200 Q/R catégorie Examens"
\end{itemize}
\end{astuce}

\subsection{Envoyer Vos Commits sur GitHub}

\begin{lstlisting}[language=bash]
# Envoyer sur GitHub
git push origin frontend
\end{lstlisting}

\textit{(Remplacez \texttt{frontend} par le nom de votre branche)}

\clearpage

% Section 7 : Créer une Pull Request
\section{Créer une Pull Request (PR)}
\label{sec:pull-request}

\subsection{Quand Créer une Pull Request ?}

\begin{itemize}
    \item[\faIcon{check}] Vous avez terminé une fonctionnalité importante
    \item[\faIcon{check}] Votre code est testé et fonctionne
    \item[\faIcon{check}] Vous voulez que votre travail soit intégré dans \texttt{dev}
\end{itemize}

\subsection{Étapes pour Créer une PR sur GitHub}

\begin{enumerate}
    \item Allez sur GitHub : \url{https://github.com/votre-compte/chatbot-supptic}
    
    \item Cliquez sur l'onglet \textbf{"Pull requests"}
    
    \item Cliquez sur \textbf{"New pull request"}
    
    \item Sélectionnez :
    \begin{itemize}
        \item \textbf{Base :} \texttt{dev} (destination)
        \item \textbf{Compare :} \texttt{frontend} (votre branche)
    \end{itemize}
    
    \item Donnez un titre clair :
    \begin{lstlisting}
Exemple : "Frontend - Interface chat complete + feedback"
    \end{lstlisting}
    
    \item Décrivez vos changements :
    \begin{lstlisting}
## Modifications
- Ajout interface chat responsive
- Integration API /chatbot/ask/
- Boutons feedback (like/dislike)
- Page statistiques basique

## Tests effectues
- Chrome, Firefox, Safari
- Mode mobile
- Connexion API testee
    \end{lstlisting}
    
    \item Cliquez sur \textbf{"Create pull request"}
\end{enumerate}

\subsection{Que Se Passe-t-il Ensuite ?}

\begin{info}[Processus de Validation]
\begin{enumerate}
    \item Les \textbf{chefs de projet} reçoivent une notification
    \item Ils examinent votre code
    \item Ils peuvent :
    \begin{itemize}
        \item \faIcon{check-circle} \textbf{Approuver} et fusionner dans \texttt{dev}
        \item \faIcon{comment} \textbf{Commenter} et demander des modifications
        \item \faIcon{times-circle} \textbf{Rejeter} si le code pose problème
    \end{itemize}
    \item Une fois approuvé, votre code est dans \texttt{dev} !
\end{enumerate}
\end{info}

\clearpage

% Section 8 : Synchronisation Quotidienne
\section{Synchronisation Quotidienne}
\label{sec:synchronisation}

\subsection{Mettre à Jour Votre Branche Chaque Matin}

\textbf{Avant de commencer à coder chaque jour :}

\begin{lstlisting}[language=bash]
# 1. Se positionner sur votre branche
git checkout frontend

# 2. Recuperer les derniers changements de dev
git pull origin dev

# 3. Fusionner dev dans votre branche
git merge dev

# 4. Resoudre les conflits eventuels (voir section suivante)

# 5. Envoyer la mise a jour
git push origin frontend
\end{lstlisting}

\begin{important}[Pourquoi Mettre à Jour ?]
D'autres équipes ont peut-être ajouté du code dans \texttt{dev} hier soir. 

En fusionnant \texttt{dev} dans votre branche chaque matin, vous évitez les conflits majeurs et restez synchronisés avec le projet.
\end{important}

\subsection{Gérer les Conflits de Fusion}

\subsubsection{Qu'est-ce qu'un Conflit ?}

Un conflit arrive quand :
\begin{itemize}
    \item Vous avez modifié un fichier
    \item Quelqu'un d'autre a modifié le \textbf{même fichier} au \textbf{même endroit}
    \item Git ne sait pas quelle version garder
\end{itemize}

\subsubsection{Résoudre un Conflit}

\begin{lstlisting}[language=bash]
# Apres git merge dev, si conflit :
# Git affiche : CONFLICT (content): Merge conflict in frontend/index.html
\end{lstlisting}

\textbf{Étapes :}

\begin{enumerate}
    \item Ouvrez le fichier en conflit dans VS Code
    
    \item Vous verrez des marqueurs :
    \begin{lstlisting}
<<<<<<< HEAD
Votre code
=======
Code de l'autre personne
>>>>>>> dev
    \end{lstlisting}
    
    \item Choisissez quelle version garder (ou combinez les deux)
    
    \item Supprimez les marqueurs \texttt{<<<<<<<}, \texttt{=======}, \texttt{>>>>>>>}
    
    \item Sauvegardez le fichier
    
    \item Ajoutez et commitez :
    \begin{lstlisting}[language=bash]
git add .
git commit -m "Resolution conflit merge dev"
git push origin frontend
    \end{lstlisting}
\end{enumerate}

\begin{astuce}[VS Code Aide Beaucoup !]
VS Code détecte automatiquement les conflits et propose des boutons :
\begin{itemize}
    \item "Accept Current Change" (garder votre version)
    \item "Accept Incoming Change" (garder leur version)
    \item "Accept Both Changes" (garder les deux)
\end{itemize}
\end{astuce}

\clearpage

% Section 9 : Commandes Git Essentielles
\section{Commandes Git Essentielles - Récapitulatif}
\label{sec:commandes}

\subsection{Configuration Initiale}

\begin{lstlisting}[language=bash]
# Cloner le depot
git clone https://github.com/compte/chatbot-supptic.git
cd chatbot-supptic

# Voir toutes les branches
git branch -a
\end{lstlisting}

\subsection{Travail Quotidien}

\begin{lstlisting}[language=bash]
# Se positionner sur votre branche
git checkout frontend

# Mettre a jour depuis dev (chaque matin)
git pull origin dev
git merge dev

# Voir les fichiers modifies
git status

# Ajouter tous les changements
git add .

# Creer un commit
git commit -m "Message clair et descriptif"

# Envoyer sur GitHub
git push origin frontend
\end{lstlisting}

\subsection{Commandes de Vérification}

\begin{lstlisting}[language=bash]
# Quelle branche suis-je ?
git branch

# Voir l'historique des commits
git log --oneline

# Voir les differences avant commit
git diff
\end{lstlisting}

\subsection{Commandes Interdites}

\begin{important}[À NE JAMAIS FAIRE]
\begin{lstlisting}[language=bash]
# INTERDIT sur main et dev
git push --force

# INTERDIT - Ne pas supprimer ces branches
git branch -D main
git branch -D dev
\end{lstlisting}
\end{important}

\clearpage

% Section 10 : Checklist Quotidienne
\section{Checklist Quotidienne}
\label{sec:checklist}

\subsection{Le Matin (Avant de Coder)}

\begin{itemize}
    \item[$\square$] Je me positionne sur ma branche : \texttt{git checkout frontend}
    \item[$\square$] Je récupère les dernières modifications : \texttt{git pull origin dev}
    \item[$\square$] Je fusionne dev dans ma branche : \texttt{git merge dev}
    \item[$\square$] Je résous les conflits éventuels
    \item[$\square$] Je vérifie que tout fonctionne (tests)
\end{itemize}

\subsection{Pendant le Travail}

\begin{itemize}
    \item[$\square$] Je code dans MON dossier uniquement
    \item[$\square$] Je teste régulièrement
    \item[$\square$] Je fais des commits réguliers (toutes les 1-2h)
    \item[$\square$] Mes messages de commit sont clairs
\end{itemize}

\subsection{Le Soir (Fin de Journée)}

\begin{itemize}
    \item[$\square$] Je vérifie que mon code fonctionne
    \item[$\square$] Je fais un dernier commit : \texttt{git add . \&\& git commit -m "..."}
    \item[$\square$] J'envoie sur GitHub : \texttt{git push origin frontend}
    \item[$\square$] Si ma feature est terminée, je crée une Pull Request
\end{itemize}

\clearpage

% Section 11 : FAQ et Problèmes Courants
\section{FAQ et Problèmes Courants}
\label{sec:faq}

\subsection{Q1 : J'ai modifié un fichier par erreur dans un autre dossier}

\textbf{Réponse :}

\begin{lstlisting}[language=bash]
# Annuler les modifications NON commitees
git checkout -- chemin/vers/fichier

# Ou annuler TOUTES les modifications non commitees
git reset --hard
\end{lstlisting}

\begin{important}[Attention]
\texttt{git reset --hard} supprime TOUTES vos modifications non sauvegardées !
\end{important}

\subsection{Q2 : J'ai fait un commit sur la mauvaise branche}

\textbf{Réponse :}

\begin{lstlisting}[language=bash]
# Annuler le dernier commit (mais garder les modifications)
git reset --soft HEAD~1

# Se positionner sur la bonne branche
git checkout ma-vraie-branche

# Refaire le commit
git add .
git commit -m "Mon message"
\end{lstlisting}

\subsection{Q3 : Mon push est rejeté par GitHub}

\textbf{Message d'erreur :}
\begin{lstlisting}
! [rejected] frontend -> frontend (non-fast-forward)
\end{lstlisting}

\textbf{Cause :} Quelqu'un a push sur la même branche pendant que vous travailliez.

\textbf{Solution :}

\begin{lstlisting}[language=bash]
# Recuperer les changements
git pull origin frontend

# Resoudre les conflits eventuels

# Re-pousser
git push origin frontend
\end{lstlisting}

\subsection{Q4 : Je ne vois pas ma branche dans VS Code}

\textbf{Solution :}

\begin{lstlisting}[language=bash]
# Mettre a jour la liste des branches
git fetch origin

# Lister toutes les branches
git branch -a

# Se positionner sur la branche
git checkout ma-branche
\end{lstlisting}

\clearpage

% Conclusion
\section*{Conclusion}

\vspace{1cm}

\begin{center}
\Large
\textbf{🎯 Résumé en 3 Points}

\vspace{0.5cm}
\end{center}

\begin{enumerate}
    \item \textbf{Travaillez sur VOTRE branche} (backend, frontend, database, data-struct)
    
    \item \textbf{Synchronisez-vous CHAQUE MATIN} avec dev
    
    \item \textbf{Proposez vos changements} via Pull Request quand c'est prêt
\end{enumerate}

\vspace{1cm}

\begin{center}

\vfill{}

\textit{Document préparé par :} \\
\textbf{NKOUMOU TJADE} \\
Chef de la Division de Programmation \\
Club Informatique SUP'PTIC

\vspace{0.5cm}

\textbf{Bonne collaboration à tous ! 🚀}
\end{center}

\end{document}